\begin{corollary}[中心六重聚合的惯性群约束：Perron 分支代数次数必被 \(6\) 整除]\label{cor:rate-center-perron-degree-multiple-6}
\leanverified{paper\_rate\_center\_perron\_degree\_multiple\_6}
沿用消元证书 \(R(\tfrac12,u)=-(u-1)^6Q(u)/64\) 的记号，并假设一般横截条件
\[
\partial_\alpha R\!\left(\tfrac12,1\right)\neq 0.
\]
令 \(F(\alpha,u)\in\QQ[\alpha,u]\) 为 \(R(\alpha,u)\) 在 \(\QQ(\alpha)[u]\) 中包含点 \((\alpha,u)=(\tfrac12,1)\) 的不可约因子（例如包含 Perron 分支的那一因子）。
则函数域扩张
\(
\QQ(\alpha)\subset \QQ(\alpha,u)\ (F(\alpha,u)=0)
\)
在 \(\alpha=\tfrac12\) 处的分歧指数为 \(e=6\)；
从而 \(F\) 的 \(u\)-次数（亦即该分支对 \(\QQ(\alpha)\) 的代数次数）必被 \(6\) 整除。
\end{corollary}
\begin{proof}
由证书 \(R(\tfrac12,u)=-(u-1)^6Q(u)/64\) 且 \(Q(1)\neq 0\)，可知在中心截面 \(\alpha=\tfrac12\) 上 \(u=1\) 为恰 \(6\) 重根。
在横截条件 \(\partial_\alpha R(\tfrac12,1)\neq 0\) 下，Weierstrass 准备定理（或代数分歧理论的等价表述）给出：在 \((\tfrac12,1)\) 的某个邻域内，方程 \(R(\alpha,u)=0\) 的局部解满足
\(
u-1\asymp (\alpha-\tfrac12)^{1/6}
\)，
即该点的局部分歧指数为 \(e=6\)（见前述 order-\(6\) singular inversion 标度）。
因此在不可约因子 \(F\) 所定义的函数域扩张中，该点对应的素位的分歧指数为 \(6\)；
而在特征 \(0\) 的代数函数域中，任一分歧指数 \(e\) 必整除扩张次数 \([\QQ(\alpha,u):\QQ(\alpha)]\)，从而 \(\deg_u F\) 必被 \(6\) 整除。证毕。
\end{proof}

\begin{remark}[Galois 闭包中的阶 \(6\) 循环惯性（语义对齐）]\label{rem:rate-center-inertia-c6}
在特征 \(0\) 下，局部惯性群为循环群，其阶等于分歧指数。
因此推论 \ref{cor:rate-center-perron-degree-multiple-6} 所述中心六重聚合在相应的 Galois 闭包中强制出现一个阶为 \(6\) 的循环惯性子群；
该约束完全由中心聚合的可审计因子 \((u-1)^6\) 强制，而不依赖额外数值拟合。
本注仅用于对齐“惯性/分歧”语义，不作为其他证明的前提。
\end{remark}
